\section{先把一次运行落实为几项人工动作}

桌面上只有一块接好电源的电路板和一台能够向存储单元写入字节的调试工具。操作人员希望
运行一段求和程序，输入是
\[
7,\quad -2,\quad 11,\quad 5,
\]
预期结果是 \(21\)。在没有操作系统的条件下，这项要求首先表现为四项具体动作：

\begin{enumerate}
\item 保持复位，使处理器暂时不能执行普通指令；
\item 把程序字节、四个输入整数和最初需要的栈内容写到约定位置；
\item 设置完成后释放复位，让处理器从固定起点开始；
\item 等待程序停止或返回，再读取结果位置中的四个字节。
\end{enumerate}

这里的调试工具只在机器停止时帮助写入 RAM。它不参与程序运行，也不是程序可以调用的
设备。采用这种安排不是因为它方便，而是因为它暴露了最朴素的方法：人在运行之前把每个
必要字节摆好，机器只负责执行。第二章才会处理这种人工装入为何难以长期使用。

“程序”这个名称暂时只表示一段按顺序解释的指令字节。“一次运行”则表示处理器从约定
起点开始，根据这些字节反复更新自身状态和存储内容，直到到达约定的结束位置。此时还没有
程序身份、进程状态或资源管理对象。

\subsection{为什么不能只有处理器和一片 RAM}

假设机器只接处理器和 RAM。处理器内部的状态在刚上电时未必具有确定值，RAM 中的字节也
会在断电后丢失。即使操作人员已经写好程序，仍有两个问题没有答案：处理器从哪个地址开始
读取；写入工具撤走以后，谁保证每次释放复位都回到相同起点。

因此，运行机器还需要一小段断电后仍保留的启动字节。它们放在非易失启动存储中。复位释放
时，处理器先从这段固定内容开始执行；固定内容检查本次人工装入的布局，准备寄存器和栈，
最后才把指令位置交给 RAM 中的求和程序。后文把这段最先执行的软件称为启动程序或固件。

RAM 仍然不可替代。求和程序要保存输入、栈和最终结果，这些内容会在运行期间改变。启动
存储负责“断电之后仍然存在”，RAM 负责“运行期间能够改写”，两者解决的是不同问题。

\subsection{五类部件怎样形成一台可运行机器}

处理器、启动存储和 RAM 还不能直接随意并联。处理器发出一个地址时，必须只有一个目标
回答；目标返回较慢时，处理器还必须等待结果。机器因而使用一组选择和转发电路，把一次
请求交给唯一器件，再把器件的响应送回处理器。后文在解释选择规则后，才把这组电路称为
地址互连。

此外，寄存器更新需要共同的时序边界，刚上电的随机状态需要被收束到固定起点。时钟源提供
周期性边沿，复位电路在电源尚未稳定时阻止普通执行，并在释放时给处理器规定初始状态。

至此，当前运行机器包含五类必要部件：

\begin{center}
\begin{tikzpicture}[
  node distance=11mm and 7mm,
  unit/.style={stage,minimum height=1.28cm},
  signal/.style={-{Stealth[length=2.2mm,width=1.6mm]},draw=BookMuted,line width=0.58pt},
  request/.style={->,>=Latex,thick,draw=BookAccent},
  response/.style={->,>=Latex,thick,draw=BookMuted},
  prep/.style={->,>=Latex,thick,dashed,draw=BookTeal}
]
\node[unit,minimum width=2.85cm] (cpu)
  {处理器\\PC、SP、寄存器堆\\指令与待处理事务};
\node[unit,right=of cpu,minimum width=3.05cm] (link)
  {地址互连\\请求锁存、范围比较\\目标选择与等待记录};
\node[unit,right=of link,minimum width=2.55cm] (owner)
  {RAM 端口选择\\复位期：夹具\\运行期：互连};
\node[unit,below=13mm of owner,minimum width=2.95cm] (ram)
  {RAM\\事务检查、字节通道\\易失字节阵列};
\node[unit,above=13mm of owner,minimum width=2.95cm] (boot)
  {启动存储\\范围检查、地址下标\\非易失字节阵列};
\node[unit,below=13mm of link,minimum width=2.75cm,fill=BookWarm] (fixture)
  {外部调试夹具\\地址/数据寄存器\\请求控制与读回};
\node[unit,above=17mm of cpu,minimum width=2.85cm] (clock)
  {时钟源\\振荡、整形与输出\\产生共同更新边沿};
\node[unit,right=of clock,minimum width=3.05cm] (reset)
  {复位电路\\电源监测、保持状态\\同步释放};
\draw[signal] (clock) -- (cpu);
\draw[signal] (clock) -- (reset);
\draw[signal] (reset.south) |- (cpu.north);
\draw[signal] (reset.south) -| (owner.north);
\draw[request] ([yshift=3.5mm]cpu.east) -- ([yshift=3.5mm]link.west);
\draw[response] ([yshift=-3.5mm]link.west) -- ([yshift=-3.5mm]cpu.east);
\draw[request] ([yshift=3mm]link.east) -| ([xshift=-4mm]boot.south);
\draw[response] ([xshift=4mm]boot.south) |- ([yshift=-3mm]link.east);
\draw[request] ([yshift=3mm]link.east) -- ([yshift=3mm]owner.west);
\draw[response] ([yshift=-3mm]owner.west) -- ([yshift=-3mm]link.east);
\draw[request] ([xshift=-4mm]owner.south) -- ([xshift=-4mm]ram.north);
\draw[response] ([xshift=4mm]ram.north) -- ([xshift=4mm]owner.south);
\draw[prep] ([xshift=4mm]fixture.north) -| ([xshift=-7mm]owner.south);
\draw[prep,<-] ([xshift=-4mm]fixture.north) -| ([xshift=7mm]owner.south);
\node[font=\scriptsize,BookMuted,above=2mm of boot] {启动地址范围};
\node[font=\scriptsize,BookMuted,below=2mm of fixture] {仅在复位保持时取得 RAM 端口};
\node[font=\scriptsize,BookAccent,below=1mm of cpu] {请求：地址、方向、宽度、写数据};
\node[font=\scriptsize,BookMuted,below=5mm of cpu] {响应：状态与读数据};
\end{tikzpicture}
\end{center}
\diagramnote{蓝色箭头表示运行期请求，灰色箭头表示响应，虚线表示复位期间的人工准备通路。
复位电路同时约束处理器与 RAM 端口选择器，因此调试夹具和处理器不会同时驱动 RAM。图中
没有操作系统、调度器、系统调用或可编程计时器。}

这幅图同时给出器件边界和端口关系，但每个框中的结构仍需分别验证。后文会依次展开处理器
如何保存指令状态，启动存储与 RAM 的阵列怎样被地址选中，互连怎样保持尚未完成的请求，
时钟与复位怎样产生确定边界，以及调试夹具怎样在人工准备结束后交还 RAM 端口。

\subsection{连接不是一条没有含义的双向线}

处理器读取时，地址和读请求从处理器流向目标，数据从目标返回。处理器写入时，地址、写
数据和允许更新的字节位置流向目标，目标返回成功或失败状态。时钟与复位则沿另一组连接
送到需要按边沿更新的部件。不同方向承载不同责任，不能用一根无说明的线概括。

\begin{longtable}{|L{0.23\linewidth}|L{0.27\linewidth}|L{0.38\linewidth}|}
\hline
\rowcolor{BookTableHead}
连接方向 & 运行时携带的内容 & 接收方必须完成的判断 \\ \hline
处理器到选择电路 & 地址、读取或写入、宽度、写数据 & 哪个器件拥有该地址 \\ \hline
选择电路到目标 & 已选中的请求和器件内部位置 & 请求是否合法，何时可以完成 \\ \hline
目标到处理器 & 成功或失败、读取数据 & 当前指令能否提交结果 \\ \hline
时钟与复位到处理器 & 周期边沿、复位是否有效 & 保持旧状态还是进入规定起点 \\ \hline
\end{longtable}

表中尚未给信号命名，因为读者首先需要理解信息方向。等到讨论快慢不同的器件时，再引入
“有效”和“就绪”等握手信号，并说明它们在哪个边沿共同成立。

\subsection{当前程序将占用哪些位置}

整体方位明确后，才需要给本次运行分配具体位置。教学处理器能够表达 32 位地址；每个地址
指出一个字节。启动存储位于低地址，RAM 位于另一段地址。求和程序、输入、结果和栈在 RAM
中彼此分开：

\begin{longtable}{|L{0.25\linewidth}|L{0.25\linewidth}|L{0.36\linewidth}|}
\hline
\rowcolor{BookTableHead}
地址或范围 & 人在释放复位前放入什么 & 运行期间由谁改变 \\ \hline
\code{0x00000000--0x0000ffff} & 固定启动程序 & 普通运行不能改写 \\ \hline
\code{0x00100000--0x0010001f} & 八条求和指令 & 本次程序不改写 \\ \hline
\code{0x00102000--0x0010200f} & 四个输入整数 & 本次程序只读取 \\ \hline
\code{0x00102020--0x00102023} & 初始清零的结果槽 & 求和程序写入 21 \\ \hline
\code{0x001efffc--0x001effff} & 返回启动程序的地址 & 程序结束时读取 \\ \hline
\end{longtable}

这些名称不是存储器内部的类型。RAM 只保存位；“指令”“输入”“结果”和“返回地址”来自
处理器在不同阶段怎样使用同一组字节。地址表目前也只是本次人工约定，尚无硬件权限阻止
程序写到别的区域。

\subsection{把目标写成可以检查的启动契约}

释放复位前，操作人员和机器之间需要一份最小约定。否则即使字节都已写入，启动程序也不
知道入口、输入数量和结果位置。当前样机采用以下具体值：

\begin{longtable}{|L{0.22\linewidth}|L{0.25\linewidth}|L{0.39\linewidth}|}
\hline
\rowcolor{BookTableHead}
约定项 & 本次取值 & 为什么必须明确 \\ \hline
求和程序入口 & \code{0x00100000} & 决定第一次任务取指的位置 \\ \hline
输入首地址 & \code{0x00102000} & 决定第一个整数从哪里读取 \\ \hline
输入数量 & \(4\) & 决定循环何时结束 \\ \hline
结果地址 & \code{0x00102020} & 让程序与观察者指向同一结果槽 \\ \hline
初始栈位置 & \code{0x001efffc} & 保存程序结束后的控制去向 \\ \hline
预期结果 & \(21\) & 提供可独立核对的外部事实 \\ \hline
\end{longtable}

启动程序会把输入首地址、数量和结果地址放进处理器寄存器，并把栈指针指向预置返回地址。
这里先说明约定的作用，寄存器和栈为何能够保存这些值将在下一节从计算中逐步得到。

\subsection{什么时候才算完成了第一次运行}

“机器开始动了”不是完成标准。对本章而言，成功至少要同时满足四项可检查事实：

\begin{enumerate}
\item 复位释放后，第一笔取指来自规定的启动位置；
\item 启动程序准备的入口、参数与栈恰好等于上表；
\item 求和程序逐条执行，结果槽最终保存字节 \code{15 00 00 00}；
\item 程序沿预置返回地址回到启动程序，启动程序停止继续改写本次结果。
\end{enumerate}

第四项把“结果已经写入”和“本次控制流已经结束”分开。程序可能先写出 21，随后又因错误
继续执行；只观察结果槽不足以证明整次运行正确。后文将分别给出结果写入和控制返回的状态
快照。

\begin{problemframe}
\textbf{本章尚未解决的问题。} 操作人员必须在每次运行前保持复位，逐区域写入程序、输入
和栈，再检查地址是否一致。机器自身既不能从外部接收一份新程序，也不能判断一串到达的
字节是否完整。先把这次人工运行讲清楚，章末再用这些具体成本推动下一次演进。
\end{problemframe}
